%%%%%%%%%%%%%%%%%%%%%%%%%%%%%%%%%%%%%%%%%%%%%%%%%%%%%%%%%%%%%%%%%%%%%%%%%%
%
% Plantilla para libro de texto de matemáticas.
%
% Esta plantilla ha sido desarrollada desde cero, pero utiliza algunas partes
% del código de la plantilla original utilizada en apuntesDGIIM
% (https://github.com/libreim/apuntesDGIIM), basada a su vez en las plantillas
% 'Short Sectioned Assignment' de Frits Wenneker (http://www.howtotex.com),
% 'Plantilla de Trabajo' de Mario Román y 'Plantilla básica de Latex en Español'
% de Andrés Herrera Poyatos (https://github.com/andreshp). También recoge
% ideas de la plantilla 'Multi-Purpose Large Font Title Page' de
% Frits Wenneker y Vel (vel@latextemplates.com).
%
% Licencia:	
% CC BY-NC-SA 4.0 (https://creativecommons.org/licenses/by-nc-sa/4.0/)
%
%%%%%%%%%%%%%%%%%%%%%%%%%%%%%%%%%%%%%%%%%%%%%%%%%%%%%%%%%%%%%%%%%%%%%%%%%

% ---------------------------------------------------------------------------
% CONFIGURACIÓN BÁSICA DEL DOCUMENTO
% ---------------------------------------------------------------------------

%\documentclass[11pt, a4paper, twoside]{article} % Usar para imprimir
\documentclass[10pt, a4paper]{article}

\usepackage[normalem]{ulem}
\usepackage{cancel}
\newcommand\dout{\bgroup \markoverwith{\rule[0.2ex]{0.1pt}{0.4pt}\rule[0.8ex]{0.1pt}{0.4pt}}\ULon}
\def\dout{\bgroup
 \markoverwith{\lower-0.35ex\hbox
 {\kern-.03em\vbox{\hrule width.2em\kern0.45ex\hrule}\kern-.03em}}%
 \ULon}
\MakeRobust\dout

\linespread{1.3}            % Espaciado entre líneas.
\setlength\parindent{0pt}   % No indentar el texto por defecto.
\setlength\parskip{7pt}

% ---------------------------------------------------------------------------
% PAQUETES BÁSICOS
% ---------------------------------------------------------------------------

% IDIOMA
\usepackage[utf8]{inputenc}
\usepackage[spanish, es-tabla, es-lcroman, es-noquoting]{babel}
\usepackage[table,xcdraw]{xcolor}

% MATEMÁTICAS
\usepackage{amsmath}    % Paquete básico de matemáticas
\usepackage{amsthm}     % Teoremas
\usepackage{mathrsfs}   % Fuente para ciertas letras utilizadas en matemáticas

% FUENTES
\usepackage{newpxtext, newpxmath}   % Fuente similar a Palatino
\usepackage{FiraSans}                 % Fuente sans serif
\usepackage[T1]{fontenc}
\usepackage[italic]{mathastext}     % Utiliza la fuente del documento
                                    % en los entornos matemáticos

% MÁRGENES
\usepackage[margin=2.5cm, top=3cm]{geometry}

% LISTAS
\usepackage{enumitem}       % Mejores listas
\setlist{leftmargin=.5in}   % Especifica la indentación para las listas.

% Listas ordenadas con números romanos (i), (ii), etc.
\newenvironment{nlist}
{\begin{enumerate}
    \renewcommand\labelenumi{(\emph{\roman{enumi})}}}
  {\end{enumerate}}

%  OTROS
\usepackage[hidelinks]{hyperref}   % Enlaces
\usepackage{graphicx}   % Permite incluir gráficos en el documento
\usepackage{relsize}

% LISTINGS
\usepackage{listings}
\usepackage{xcolor}     % Permite definir y utilizar colores
\usepackage{lipsum}
\usepackage{courier}

% Fijar tabla a posición
\usepackage{array}
\newcolumntype{L}[1]{>{\raggedright\let\newline\\\arraybackslash\hspace{0pt}}m{#1}}
\newcolumntype{C}[1]{>{\centering\let\newline\\\arraybackslash\hspace{0pt}}m{#1}}
\newcolumntype{R}[1]{>{\raggedleft\let\newline\\\arraybackslash\hspace{0pt}}m{#1}}


% Colores para los bloques de código
\definecolor{codegreen}{rgb}{0,0.6,0}
\definecolor{codegray}{rgb}{0.5,0.5,0.5}
\definecolor{codepurple}{rgb}{0.58,0,0.82}
\definecolor{backcolour}{rgb}{0.95,0.95,0.92}
\lstdefinestyle{mystyle}{
	backgroundcolor=\color{backcolour},
	commentstyle=\color{codegreen},
	keywordstyle=\color{blue},
	numberstyle=\tiny\color{codegray},
	stringstyle=\color{codepurple},
	basicstyle=\footnotesize\ttfamily,
	breakatwhitespace=false,         
	breaklines=true,                 
	captionpos=b,                    
	keepspaces=true,                 
	numbers=left,                    
	numbersep=5pt,                  
	showspaces=false,                
	showstringspaces=false,
	showtabs=false,                  
	tabsize=4
}
\lstset{style=mystyle}

%\lstset{basicstyle=\footnotesize\ttfamily,breaklines=true}
%\lstset{framextopmargin=50pt,frame=bottomline}
 
% ---------------------------------------------------------------------------
% COMANDOS PERSONALIZADOS
% ---------------------------------------------------------------------------

% \equalto
\newcommand{\verteq}{\rotatebox{90}{$\,=$}}
\newcommand{\equalto}[2]{\underset{\scriptstyle\overset{\mkern4mu\verteq}{#2}}{#1}}


% ---------------------------------------------------------------------------
% COLORES
% ---------------------------------------------------------------------------

\definecolor{50}{HTML}{e8f5e9}
\definecolor{100}{HTML}{c8e6c9}
\definecolor{200}{HTML}{a5d6a7}
\definecolor{300}{HTML}{81c784}
\definecolor{400}{HTML}{66bb6a}
\definecolor{500}{HTML}{4caf50}
\definecolor{600}{HTML}{43a047}
\definecolor{700}{HTML}{388e3c}
\definecolor{800}{HTML}{2e7d32}
\definecolor{900}{HTML}{1b5e20}

% ---------------------------------------------------------------------------
% DISEÑO DE PÁGINA
% ---------------------------------------------------------------------------

\usepackage{pagecolor}
\usepackage{afterpage}

% ---------------------------------------------------------------------------
% CABECERA Y PIE DE PÁGINA
% ---------------------------------------------------------------------------

\usepackage{fancyhdr}   % Paquete para cabeceras y pies de página

% Indica que las páginas usarán la configuración de fancyhdr
\pagestyle{fancy}
\fancyhf{}

% Representa la sección de la cabecera
\renewcommand{\sectionmark}[1]{%
\markboth{#1}{}}

% Representa la subsección de la cabecera
\renewcommand{\subsectionmark}[1]{%
\markright{#1}{}}

% Parte derecha de la cabecera
\fancyhead[LE,RO]{\sffamily\textsl{\rightmark} \hspace{1em}  \textcolor{500}{\rule[-0.4ex]{0.2ex}{1.2em}} \hspace{1em} \thepage}

% Parte izquierda de la cabecera
\fancyhead[RE,LO]{\sffamily{}}

% Elimina la línea de la cabecera
\renewcommand{\headrulewidth}{0pt}

% Controla la altura de la cabecera para que no haya errores
\setlength{\headheight}{14pt}

% ---------------------------------------------------------------------------
% TÍTULOS DE PARTES Y SECCIONES
% ---------------------------------------------------------------------------


\usepackage{titlesec}

% Estilo de los títulos de las partes
\titleformat{\part}[hang]{\Huge\bfseries\sffamily}{\thepart\hspace{20pt}\textcolor{500}{|}\hspace{20pt}}{0pt}{\Huge\bfseries}
\titlespacing*{\part}{0cm}{-2em}{2em}[0pt]


\usepackage{listingsutf8}
% Reiniciamos el contador de secciones entre partes (opcional)
\makeatletter
\@addtoreset{section}{part}
\makeatother

% Estilo de los títulos de las secciones, subsecciones y subsubsecciones
\titleformat{\section}
  {\Large\bfseries\sffamily}{\thesection}{0em}{}

\titleformat{\subsection}
  {\Large\sffamily}{\thesubsection}{1em}{}[\vspace{.5em}]

\titleformat{\subsubsection}
  {\sffamily}{\thesubsubsection}{1em}{}

% CAMBIO DE NUMERACIÓN - ESPECÍFICO DE ESTE DOCUMENTO
\renewcommand{\thesection}{}
\renewcommand{\thesubsection}{\arabic{subsection}}

% ---------------------------------------------------------------------------
% ENTORNOS PERSONALIZADOS
% ---------------------------------------------------------------------------

\usepackage{mdframed}

%% DEFINICIONES DE LOS ESTILOS

% Nuevo estilo para definiciones
\newtheoremstyle{definition-style}  % Nombre del estilo
{}                                  % Espacio por encima
{}                                  % Espacio por debajo
{}                                  % Fuente del cuerpo
{}                                  % Identación
{\bf\sffamily}                      % Fuente para la cabecera
{.}                                 % Puntuación tras la cabecera
{.5em}                              % Espacio tras la cabecera
{\thmname{#1}\thmnumber{ #2}\thmnote{ (#3)}}  % Especificación de la cabecera

% Nuevo estilo para notas
\newtheoremstyle{remark-style}
{10pt}
{10pt}
{}
{}
{\itshape \sffamily}
{.}
{.5em}
{}

% Nuevo estilo para teoremas y proposiciones
\newtheoremstyle{theorem-style}
{}
{}
{}
{}
{\bfseries \sffamily}
{.}
{.5em}
{\thmname{#1}\thmnumber{ #2}\thmnote{ (#3)}}

% Nuevo estilo para ejemplos
\newtheoremstyle{example-style}
{10pt}
{10pt}
{}
{}
{\bf \sffamily}
{}
{.5em}
{\thmname{#1}\thmnumber{ #2.}\thmnote{ #3.}}

% Nuevo estilo para la demostración

\makeatletter
\renewenvironment{proof}[1][\proofname] {\par\pushQED{\qed}\normalfont\topsep6\p@\@plus6\p@\relax\trivlist\item[\hskip\labelsep\itshape\sffamily#1\@addpunct{.}]\ignorespaces}{\popQED\endtrivlist\@endpefalse}
\makeatother

%% ASIGNACIÓN DE LOS ESTILOS

% Teoremas, proposiciones y corolarios
\newtheoremstyle{theorem-style}{}{}{}{}{}{}{ }{}
\theoremstyle{theorem-style}
\newtheorem*{datos}{}
\theoremstyle{theorem-style}
\newtheorem{nth}{Teorema}[section]
\newtheorem{nprop}{Proposición}[section]
\newtheorem{ncor}{Corolario}[section]
\newtheorem{lema}{Lema}[section]

% Definiciones
\theoremstyle{definition-style}
\newtheorem{ndef}{Definición}[section]

% Notas
\theoremstyle{remark-style}
\newtheorem*{nota}{Nota}

% Ejemplos
\theoremstyle{example-style}
\newtheorem{ejemplo}{Ejemplo}[section]

% Ejercicios y solución
\theoremstyle{definition-style}
\newtheorem{ejer}{Ejercicio}[section]

\theoremstyle{remark-style}
\newtheorem*{sol}{Solución}

%% MARCOS DE LOS ESTILOS

% Configuración general de mdframe, los estilos de los teoremas, etc
\mdfsetup{
  skipabove=1em,
  skipbelow=1em,
  innertopmargin=1em,
  innerbottommargin=1em,
  splittopskip=2\topsep,
}

% Definimos los marcos de los estilos


\mdfdefinestyle{datos-frame}{
	linewidth=2pt, %
	linecolor= 500, %
	topline=false, %
	bottomline=false, %
	rightline=false,%
	leftmargin=0em, %
	innerleftmargin=1em, %
	innerrightmargin=1em,
	rightmargin=0em, %
}%
\mdfdefinestyle{nth-frame}{
	linewidth=2pt, %
	linecolor= 500, %
	topline=false, %
	bottomline=false, %
	rightline=false,%
	leftmargin=0em, %
	innerleftmargin=1em, %
  innerrightmargin=1em,
	rightmargin=0em, %
}%

\mdfdefinestyle{nprop-frame}{
	linewidth=2pt, %
	linecolor= 300, %
	topline=false, %
	bottomline=false, %
	rightline=false,%
	leftmargin=0pt, %
	innerleftmargin=1em, %
	innerrightmargin=1em,
	rightmargin=0pt, %
}%

\mdfdefinestyle{ndef-frame}{
	linewidth=2pt, %
	linecolor= 500, %
	backgroundcolor= 50,
	topline=false, %
	bottomline=false, %
	rightline=false,%
	leftmargin=0pt, %
	innerleftmargin=1em, %
	innerrightmargin=1em,
	rightmargin=0pt, %
}%

\mdfdefinestyle{ejer-frame}{
	linewidth=2pt, %
	linecolor= 300, %
	backgroundcolor= 50,
	topline=false, %
	bottomline=false, %
	rightline=false,%
	leftmargin=0pt, %
	innerleftmargin=1em, %
	innerrightmargin=1em,
	rightmargin=0pt, %
}%

\mdfdefinestyle{ejemplo-frame}{
	linewidth=0pt, %
	linecolor= 300, %
	leftline=false, %
	rightline=false, %
	leftmargin=0pt, %
	innerleftmargin=1.3em, %
	innerrightmargin=1em,
	rightmargin=0pt, %
	innertopmargin=0em,%
	innerbottommargin=0em, %
	splittopskip=\topskip, %
}%

% Asignamos los marcos a los estilos
\surroundwithmdframed[style=nth-frame]{nth}
\surroundwithmdframed[style=datos-frame]{datos}
\surroundwithmdframed[style=nprop-frame]{nprop}
\surroundwithmdframed[style=nprop-frame]{ncor}
\surroundwithmdframed[style=ndef-frame]{ndef}
\surroundwithmdframed[style=ejer-frame]{ejer}
\surroundwithmdframed[style=ejemplo-frame]{ejemplo}
\surroundwithmdframed[style=ejemplo-frame]{sol}

% ---------------------------------------------------------------------------
% CONFIGURACIÓN DE LA PORTADA
% ---------------------------------------------------------------------------

\newcommand{\asignatura}{Análisis de eficiencia de algoritmos}

\newcommand{\autor}{Miguel Ángel Fernández Gutiérrez\\Pedro Gallego López\\\hspace{1cm}}

\newcommand{\grado}{Doble Grado en Ingeniería Informática y Matemáticas}

\newcommand{\universidad}{Universidad de Granada}

\newcommand{\enlaceweb}{github.com/DGIIMUnderground}

% ---------------------------------------------------------------------------
% CONFIGURACIÓN PERSONALIZADA
% ---------------------------------------------------------------------------

%%%%%%%%%%%%%%%%%%%%%%%%%%%%%%%%%%%%%%%%%%%%%%%%%%%%%%%%%%%%%%%%%%%%%%%%%%%%%
% ---------------------------------------------------------------------------
% COMIENZO DEL DOCUMENTO
% ---------------------------------------------------------------------------
%%%%%%%%%%%%%%%%%%%%%%%%%%%%%%%%%%%%%%%%%%%%%%%%%%%%%%%%%%%%%%%%%%%%%%%%%%%%%

\begin{document}

% ---------------------------------------------------------------------------
% PORTADA EXTERIOR
% ---------------------------------------------------------------------------

\newpagecolor{500}\afterpage{\restorepagecolor} % Color de la página
\begin{titlepage}

  % Título del documento
	\parbox[t]{\textwidth}{
			\raggedright % Texto alineado a la izquierda
			\fontsize{40pt}{40pt}\selectfont\sffamily\color{white}{
				\textbf{\Huge{FR | Prácticas}}\\\textbf{CarriCar}\\[0.4cm]\huge{Definición e implementación de un\\[-0.8cm]protocolo de aplicación en cliente-servidor}
      }
	}

	\vfill

	%% Autor e información del documento
	\parbox[t]{\textwidth}{
		\raggedright % Texto alineado a la izquierda
		\sffamily\large\color{white}
		\grado\\
		{\Large \autor }\\[15pt]
		\includegraphics[width=130pt]{ugrlogo.pdf}
	}

\end{titlepage}

% ---------------------------------------------------------------------------
% PÁGINA DE LICENCIA
% ---------------------------------------------------------------------------

\thispagestyle{empty}
\null
\vfill

%% Información sobre la licencia
\parbox[t]{\textwidth}{
  \includegraphics{by-nc-sa.pdf}\\[4pt]
  \raggedright % Texto alineado a la izquierda
  \sffamily\large
  {\Large Este trabajo se distribuye bajo una licencia CC BY-NC-SA 4.0.}\\[4pt]
  Eres libre de distribuir y adaptar el material siempre que reconozcas a los\\
  autores originales del documento, no lo utilices para fines comerciales\\
  y lo distribuyas bajo la misma licencia.\\[4pt]
  \texttt{creativecommons.org/licenses/by-nc-sa/4.0/}
}

% ---------------------------------------------------------------------------
% PORTADA INTERIOR
% ---------------------------------------------------------------------------

\begin{titlepage}

  % Título del documento
  \parbox[t]{\textwidth}{
  	\raggedright % Texto alineado a la izquierda
  	\fontsize{40pt}{40pt}\selectfont\sffamily\color{500}{
  		\textbf{\Huge{FR | Prácticas}}\\\textbf{CarriCar}\\[0.4cm]\huge{Definición e implementación de un\\[-0.8cm]protocolo de aplicación en cliente-servidor}
  	}
  }

	\vfill
	
	%% Autor e información del documento
	\parbox[t]{\textwidth}{
		\raggedright % Texto alineado a la izquierda
		\sffamily\large
		\grado\\
		{\Large \autor }\\[15pt]
		\includegraphics[width=130pt]{ugrlogo-dark.pdf}
	}

\end{titlepage}

% ---------------------------------------------------------------------------
% ÍNDICE
% ---------------------------------------------------------------------------

\thispagestyle{empty}
\tableofcontents
\newpage

% ---------------------------------------------------------------------------
% CONTENIDO
% ---------------------------------------------------------------------------

\part{¿Qué es una criptomoneda?}\label{quuxe9-es-una-criptomoneda}

\textbf{Bitcoin} es una de las criptomonedas más conocidas. Sepamos bien
cómo funciona o no, todos hemos oído hablar de Bitcoin. Sin embargo, hay
muchas otras, que se basan en los mismos principios básicos. Por tanto,
para explicar las criptomonedas, veamos cuáles son estos principios.

Supongamos que queremos crear una criptomoneda, \emph{DGIIMCoin}. La
primera pregunta sería para qué, y la segunda --en caso de que el para
qué nos haya convencido-- es qué necesitamos para ello.

Crear todo esto tendría sentido si queremos crear un sistema anónimo,
descentralizado y seguro para intercambiar dinero o información. Veamos
bien qué quiere decir esto:

\begin{itemize}
\itemsep1pt\parskip0pt\parsep0pt
\item
  Anónimo: los usuarios deben ser identificables para el sistema, pero
  no
\item
  Descentralizado: los sistemas convencionales depositan la confianza de
  instituciones en terceros --un ejemplo de esto sería Amazon en el
  envío de bienes, un banco, etc.--, en el caso de las criptomonedas su
  estado es mantenido con un consenso entre todos los usuarios.
\item
  Seguro:
\end{itemize}

Bien, para todo esto nuestro sistema debería tener una serie de
requisitos:

\begin{enumerate}
\def\labelenumi{\arabic{enumi}.}
\itemsep1pt\parskip0pt\parsep0pt
\item
  El sistema no requiere de una \textbf{autoridad central}: su estado es
  mantenido mediante un \textbf{consenso} distribuido.
\item
  El sistema mantiene un \textbf{seguimiento} de todas las unidades de
  la criptomoneda y de sus propietarios.
\item
  El sistema define si se pueden crear \textbf{nuevas unidades} de la
  criptomoneda. Si esto es así, el sistema debe definir las
  circunstancias de su origen y cómo determinar quién será el
  propietario de éstas.
\item
  La propiedad de las monedas puede probarse de forma exclusivamente
  \textbf{criptográfica}.
\item
  El sistema permite la realización de \textbf{transacciones} de forma
  controlada.
\end{enumerate}

Nuestro sistema girará en torno a las siguientes ideas clave:

\begin{itemize}
\itemsep1pt\parskip0pt\parsep0pt
\item
  \textbf{Descentralización.}
\item
  \textbf{Firmas digitales.}
\item
  \textbf{La contabilidad (\emph{ledger}) es la propia moneda.}
\item
  \textbf{Mecanismo de consenso.}
\item
  \textbf{Blockchain.}
\end{itemize}

\section{Compartiendo información de forma descentralizada: redes
P2P}\label{compartiendo-informaciuxf3n-de-forma-descentralizada-redes-p2p}

Para poder cumplir los dos primeros requisitos de nuestra lista,
eliminando así autoridades centralizadas, temos ya una solución
disponible: las \textbf{redes \emph{peer-to-peer}}, o \textbf{redes
P2P}.

La idea principal de estas redes es que la información se comparta de
igual a igual, al contrario a las redes cliente-servidor. En este caso,
todos los ordenadores tienen el mismo peso.

{[}Más de P2P{]}

En definitiva, podemos usar los protocolos P2P existentes para
implementar \emph{DGIIMCoin}. Genial, continuemos.

\section{Firmas digitales, criptografía y
\emph{hashing}}\label{firmas-digitales-criptografuxeda-y-hashing}

Para que nuestro sistema funcione son esenciales las \textbf{identidades
digitales}, es decir, formas de verificar que una transacción se ha
hecho realmente (alguien de verdad ha enviado dinero a otra persona).
Para eso, usamos la \textbf{criptografía}.

{[}Más de criptografía y hashing{]}

\section{Enviando transacciones a la
red}\label{enviando-transacciones-a-la-red}

Ya casi estamos. Hemos implementado comunicación P2P, mecanismos para
crear identidades digitales y formas para que los usuarios puedan firmar
y garantizar que la información es correcta. Ahora sólo nos queda enviar
información al sistema.

Como ya sabemos, no tenemos una autoridad central que valide cuánto
dinero tenemos, pero tampoco hace falta: aquí reside una de las ideas,
que la \textbf{contabilidad es la propia moneda}. Esto se debe a que
para saber cuánto dinero tenemos, simplemente nos basta con tener una
lista de todas las transacciones que hemos efectuado. Supongamos que tu
historial de transacciones contiene la información:

\begin{enumerate}
\def\labelenumi{\arabic{enumi}.}
\itemsep1pt\parskip0pt\parsep0pt
\item
  Tengo 500 \dout{D}.
\item
  Envío 20 \dout{D} a alguien para unos apuntes de Modelos de
  Computación (incluiremos su clave pública).
\item
  Quiero enviar 1 \dout{D} como impuesto de transacción al sistema (lo
  veremos más adelante).
\item
  Quiero quedarme con las \emph{DGIIMCoins} restantes.
\end{enumerate}

Lo único que queda es usar la red P2P para enviar estas transacciones al
resto de usuarios del sistema, sin olvidar de firmarlas usando nuestra
clave privada. Tras esto, todo el mundo sabrá que has pagado 20 \dout{D}
a tu compañero por apuntes.

Listo. Sin embargo, no tendrás los apuntes hasta que la red coincida en
que inicialmente tenías 500 \dout{D}, y por tanto dicha transacción
puede efectuarse. Una vez que la transacción se valide, tu compañero te
dará los apuntes.

\section{Mecanismos de verificación de
transacciones}\label{mecanismos-de-verificaciuxf3n-de-transacciones}

\subsection{Mining: proof-of-work}\label{mining-proof-of-work}

\subsection{Otros mecanismos}\label{otros-mecanismos}

\subsubsection{Proof-of-burn}\label{proof-of-burn}

\subsubsection{Proof-of-stake}\label{proof-of-stake}

\subsubsection{Proof-of-capacity}\label{proof-of-capacity}

\subsubsection{Proof-of-elapsed time}\label{proof-of-elapsed-time}

\subsubsection{Proof-of-importance}\label{proof-of-importance}

\subsection{El blockchain}\label{el-blockchain}

\subsection{Controlando la fuente de
dinero}\label{controlando-la-fuente-de-dinero}

\part{Aplicaciones de las
criptomonedas}\label{aplicaciones-de-las-criptomonedas}

\subsubsection{Más allá del dinero: criptomonedas de segunda y tercera
generación}\label{muxe1s-alluxe1-del-dinero-criptomonedas-de-segunda-y-tercera-generaciuxf3n}

\subsubsection{Normalización de las
criptomonedas}\label{normalizaciuxf3n-de-las-criptomonedas}

\subsection{Nuevas criptomonedas}\label{nuevas-criptomonedas}

Criptomonedas de segunda y tercera generación. Conceptos.

\subsection{Demo: Wireshark}\label{demo-wireshark}

\dout{D}

\end{document}
 